% ===========================================================================
%  Apêndices
% ===========================================================================
\chapter{Ficha de Exemplo --- Roxy Migurdia}

Use esta ficha pra calibrar se os seus números parecem certos na mesa.

\textbf{Roxy Migurdia} --- Migurd, Santo de Água, Avançado de Terra e Vento, Intermediário
de Cura.

\begin{itemize}
    \item \textbf{Atributos:} Força 0 \textbullet\ Agilidade 3 \textbullet\ Vigor 2
          \textbullet\ Intelecto 6 (já com $+1$ de Migurd) \textbullet\ Espírito 5
    \item \textbf{PV:} corpo treinado ($20 + $ o dobro dos dados de PV dos 12 ranks dela,
          nas 4 árvores $= 20 + 142 = 162$) $\times$ Fator de Vigor 2 ($\times$1,40) $=$
          \textbf{226 PV}
    \item \textbf{PM:} só a melhor escola conta, nunca a soma de todas --- o maior entre
          Espírito 5 e 4, vezes o Bônus do Santo de Água (4), $+8 = 28$, mais os $+6$ PM
          fixos da raça Migurd $=$ \textbf{34 PM}
    \item \textbf{BC de Água:} $6 + 4 = 10$ $\to$ acerta com 1d20$+$10, CD 18, dano $+10$
    \item \textbf{BC de Cura:} Espírito $5 + 2 = 7$ --- ela fecha ferimento, mas o
          Intermediário dela não salva ninguém de uma ferida mortal
    \item \textbf{CA:} 13
    \item \textbf{PP e PT:} nenhum. Ela não tem patamar em árvore do Corpo nem de Utilidade
\end{itemize}

\begin{nota}{Leitura da ficha}
    Ela acerta praticamente qualquer coisa, tem uma reserva de mana que sustenta um combate
    longo inteiro, e cai em poucos golpes de qualquer espadachim decente --- 226 PV é bastante
    numa conta isolada, mas fica baixo perto de um personagem do Corpo com a mesma quantidade
    de Ranks investidos, cujos dados de PV por patamar são bem maiores. É exatamente isso que
    ela é na história: uma professora genial dentro de um corpo frágil, que sobrevive porque
    nunca deixa ninguém chegar perto.

    \medskip
    Repare no que o Fator de Vigor faz aqui: com Vigor 2 ela multiplica por 1,40. Se tivesse
    largado Vigor em $-2$ pra comprar mais um ponto de Intelecto, o mesmo corpo treinado de
    162 viraria \textbf{64 PV} --- e um único golpe de espadachim Santo resolveria a luta.
\end{nota}

% ---------------------------------------------------------------------------
\chapter{Molde para Novas Escolas}

Cada escola nova precisa exatamente destes oito itens:

\begin{enumerate}
    \item Uma \textbf{condição-assinatura} que a escola aplica de graça (Água $\to$ Molhado).
    \item Um \textbf{combo interno} que paga por aplicar a condição (Água: gelo dobra frio
          contra Molhado; eletricidade dobra tudo).
    \item Uma \textbf{curva de PV/PM própria} que diferencie a escola. Água: PM alto, PV
          médio. Terra: PV alto, PM baixo. Fogo: dano alto, defesa nenhuma. Vento:
          meio-termo com bônus de deslocamento.
    \item \textbf{Seis Maestrias automáticas}, uma por rank --- a do Avançado sempre destranca
          Magia Combinada, a do Rei sempre destranca um elemento secundário (Água $\to$
          Eletricidade, Fogo $\to$ Explosão/Plasma, Vento $\to$ Som/Vácuo, Terra $\to$
          Metal/Magma).
    \item Uma \textbf{Magia Assinatura \assinatura} por rank, custando $+1$ PA.
    \item Uma \textbf{magia de utilidade pura} que não causa dano nenhum, mas define a
          identidade da escola fora de combate.
    \item De \textbf{6 a 8 conhecimentos} por rank baixo, 3 a 4 por rank alto --- o suficiente
          pra tabela de desbloqueio fechar sem obrigar o jogador a comprar magia velha só pra
          bater a contagem.
    \item Declare \textbf{qual atributo alimenta o BC} da escola --- Fogo, Água, Vento e Terra
          usam Intelecto; Cura, Barreira, Desintoxicação e Invocação usam Espírito. Isso
          divide a Árvore da Magia em duas metades que não competem pelos mesmos pontos.
\end{enumerate}

% ---------------------------------------------------------------------------
\chapter{Tabela Comparativa de Dano por Turno}

A régua com que toda árvore futura deve ser medida. Valores médios, alvo de CA razoável,
atributo principal progredindo de 4 até 8.

\begin{tabelalivro}{C{1.2cm} C{1.9cm} C{1.3cm} C{1.3cm} C{1.3cm} C{1.2cm} C{1.5cm} C{1.5cm} C{1.6cm}}{Árvore da Magia}
\cabtabela{\textbf{Pat.} & \textbf{Água} & \textbf{Fogo} & \textbf{Vento} & \textbf{Terra} & \textbf{Cura} & \textbf{Desint.} & \textbf{Barreira} & \textbf{Invoc.}}
1º & \textasciitilde10 & \textasciitilde12 & \textasciitilde9 & \textasciitilde11 & --- & --- & --- & \textasciitilde13 \\
2º & \textasciitilde20 & \textasciitilde26 & \textasciitilde18 & \textasciitilde24 & --- & \textasciitilde8 & --- & \textasciitilde24 \\
3º & \textasciitilde28 & \textasciitilde40 & \textasciitilde30 & \textasciitilde36 & --- & \textasciitilde20 & --- & \textasciitilde38 \\
4º & \textasciitilde22 $+$ área & \textasciitilde62 & \textasciitilde45 & \textasciitilde52 & --- & \textasciitilde28 & --- & \textasciitilde55 \\
5º & \textasciitilde54 & \textasciitilde90 & \textasciitilde70 & \textasciitilde76 & \textasciitilde40 & \textasciitilde55 & \textasciitilde30 & \textasciitilde80 \\
6º & \textasciitilde39 em 45m & \textasciitilde130 & \textasciitilde110 & \textasciitilde105 & \textasciitilde55 & \textasciitilde70 & \textasciitilde40 & \textasciitilde110 \\
\end{tabelalivro}

\begin{tabelalivro}{C{1.2cm} C{1.7cm} C{1.5cm} C{1.6cm} C{1.4cm} C{1.6cm} C{1.6cm} C{2.4cm}}{Corpo e Utilidade}
\cabtabela{\textbf{Pat.} & \textbf{Espada} & \textbf{Norte} & \textbf{Suishin} & \textbf{Arco} & \textbf{Lutador} & \textbf{Escudos} & \textbf{Utilidade}}
1º & \textasciitilde25 & \textasciitilde19 & \textasciitilde11 & \textasciitilde22 & \textasciitilde21 & \textasciitilde7 & \textasciitilde15 (1º turno) \\
2º & \textasciitilde34 & \textasciitilde25 & \textasciitilde26 & \textasciitilde34 & \textasciitilde32 & \textasciitilde9 & \textasciitilde18 \\
3º & \textasciitilde62 & \textasciitilde34 & \textasciitilde40 & \textasciitilde48 & \textasciitilde44 & \textasciitilde11 & \textasciitilde22 \\
4º & \textasciitilde78 & \textasciitilde42 & \textasciitilde60 & \textasciitilde62 & \textasciitilde58 & \textasciitilde13 & \textasciitilde26 \\
5º & \textasciitilde98 & \textasciitilde55 & \textasciitilde85 & \textasciitilde78 & \textasciitilde74 & \textasciitilde15 & \textasciitilde30 \\
6º & \textasciitilde118 & \textasciitilde87 & 0 a $\infty$ & \textasciitilde91 & \textasciitilde95 & \textasciitilde18 & \textasciitilde31 \\
\end{tabelalivro}

\begin{aviso}{Três coisas que a tabela não diz sozinha}
    \textbf{A Espada conta 4 Ações do Avançado em diante.} A Maestria \enquote{Velocidade
    Encarnada} dá uma Ação extra a quem não se move no turno, e os números já assumem isso.
    Ela é a única coluna com uma 4ª Ação antes do Imperador.

    \medskip
    \textbf{Escudos pressupõe todas as Ações gastas defendendo.} Um Defensor Imperador que
    \emph{escolha} atacar faz perto de 48 por turno, não 18. A coluna mede o que ele faz no
    papel dele, não o teto dele.

    \medskip
    \textbf{Magia não está amortizada pelas Ações.} Uma magia de Imperador custa 6 Ações ---
    dois turnos inteiros. O Sol Menor aparece como \textasciitilde130, mas entrega
    \textasciitilde65 por turno. Compare marcial com marcial e magia com magia; cruzar as
    duas metades desta tabela engana.
\end{aviso}

\begin{nota}{Como ler esta tabela}
    Número alto não significa personagem melhor. Significa personagem \emph{mais estreito}.

    \begin{itemize}
        \item O \textbf{Fogo} tem o maior número e o menor corpo --- 33 de soma de dados de
              PV, o menor do livro. Mata tudo, morre de qualquer coisa, e queima o saque.
        \item A \textbf{Água} tem o menor número entre as ofensivas e vence campanhas --- o
              valor dela é em área, a 45 metros, com aliados poupados e sem chance de errar.
        \item A \textbf{Terra} é a única que constrói. Metade do valor dela nunca aparece
              aqui: pontes, fortalezas, masmorras vedadas.
        \item O \textbf{Arco} só é real contra quem não veste aura. Contra um Santo ou
              superior, subtraia o dobro do Bônus de Rank do alvo de cada disparo.
        \item O \textbf{Suishin-ryū} não tem número. Contra quatro inimigos agressivos ele bate
              mais que qualquer coisa deste livro. Contra um inimigo parado, causa zero.
        \item O \textbf{Lutador} tem o número errado na tabela --- o que ele realmente faz é
              acumular Quebrantado. No quarto turno o inimigo já perdeu 6 de CA e 6 de dano.
        \item \textbf{Escudos} é a menor coluna do livro e o personagem mais difícil de
              substituir. Não causa dano. Decide quem sobrevive.
        \item \textbf{Cura, Desintoxicação e Barreira} não deveriam estar nesta tabela ---
              estão só pra deixar claro que, se você escolher uma delas esperando causar dano,
              escolheu errado.
    \end{itemize}
\end{nota}

% ---------------------------------------------------------------------------
\chapter{Ambiguidades Resolvidas}

Perguntas que a mesa vai fazer, respondidas antes de virarem discussão.

\subsection{Sobre Ranks e Múltiplas Árvores}

\begin{description}
    \item[Tenho Norte Santo e Espada Principiante. Faço um ataque comum com a espada.
          Quantos degraus?] Um ataque comum usa os degraus do seu maior patamar entre as
          árvores do Corpo --- os quatro do Norte. Uma técnica usa sempre os degraus da árvore
          que a concedeu: a Espada de Luz do seu Principiante rola com um degrau só, e é por
          isso que ela é ruim na sua mão.
    \item[Um talento diz \enquote{seu Bônus de Rank} e eu tenho cinco árvores. Qual uso?] O da
          árvore que concedeu o talento. Se a regra for genérica do livro e não citar árvore,
          use o maior que você tiver.
    \item[Em qual perícia meu Bônus de Rank de Utilidade soma?] Apenas nas que aquela árvore
          cobre (Cap. 3). Se duas árvores cobrem a mesma perícia, use o maior --- não some.
\end{description}

\subsection{Sobre Touki e PT}

\begin{description}
    \item[Abri o Deus da Espada no 2º patamar. Ele conta pro $+1$ por patamar?] Sim --- única
          exceção do livro. O Deus da Espada conta o 2º patamar dele nessa soma, porque foi lá
          que a aura acordou.
    \item[Sou mago de Terra Imperador. Tenho PT?] Não. Nenhum PT, em patamar nenhum. Magia e
          Utilidade nunca recebem Touki.
    \item[O Manto de Touki funciona em Postura de Água?] Funciona --- são coisas diferentes.
          Os bônus de CA das duas não somam (regra de empilhamento); use o maior.
    \item[Perdi todos os PT. Perco o Manto?] Não. O Manto é gratuito e passivo. PT pagam
          manobras, não a aura.
    \item[\enquote{Aguentar} é a manobra de 2 PT ou a técnica de Escudos?] A manobra. A técnica
          de Cavalaria e Escudos chama-se \textbf{Aguentar o Baque} --- nome diferente
          justamente porque um personagem de Escudos tem as duas.
\end{description}

\subsection{Sobre Cura, Aflições e Descanso}

\begin{description}
    \item[Magia de Cura cura veneno?] Não. Nunca, em patamar nenhum. Isso é Desintoxicação, e
          a separação é absoluta.
    \item[Desintoxicação cura PV?] Não. Ela remove a causa; a carne continua aberta.
    \item[Ferida Selada conta como Ferida Fresca pra cura em dobro?] Conta. É exatamente pra
          isso que Selar a Ferida existe.
    \item[Quantos Descansos Curtos por dia?] Dois, entre dois Descansos Longos. PM voltam pela
          metade; PT voltam inteiros; PV e PP não voltam.
\end{description}

\subsection{Sobre Condições}

\begin{description}
    \item[Posso estar Molhado e Em Chamas ao mesmo tempo?] Não. Fogo em alvo Molhado evapora a
          água (a condição some, o alvo sofre $+2$ pelo choque térmico, e não pega fogo naquele
          golpe). Água em alvo Em Chamas apaga o fogo e aplica Molhado.
    \item[Quebrantado some com magia de Cura?] Não. Não é ferimento --- é o corpo parando de
          responder. Só um Descanso Curto limpa.
    \item[Desequilibrado tira todas as minhas Reações?] Não. Limita a uma por rodada.
    \item[Congelado e Atolado juntos deixam o Deslocamento negativo?] Deslocamento não fica
          abaixo de 0. As condições não se somam em efeito, mas escapar exige resolver as duas.
    \item[O que exatamente Selado impede?] Só magia acima do rank em Barreira de quem ergueu a
          barreira. Touki, armas e habilidades de Utilidade funcionam normalmente lá dentro.
\end{description}

\subsection{Sobre Ações e Reações}

\begin{description}
    \item[Quantas Reações eu tenho?] Uma por rodada, sempre --- a menos que um efeito diga o
          contrário.
    \item[Posso usar a Reação no meu próprio turno?] Pode, desde que o gatilho aconteça.
    \item[Conjurar uma magia de 4 Ações me deixa sem Reação?] Não. Reação é independente do
          custo em Ações --- mas se sofrer dano, faça o teste de Interrupção (CD 10 $+$ o Bônus
          de Rank de quem acertou) ou perde o cântico.
    \item[O invocado gasta minhas Ações?] Ordens gerais, não. Ordens específicas, 1 Ação sua.
          Ele tem Iniciativa própria e age sozinho no turno dele.
    \item[A Silenciosa gratuita de Principiante conta no Teto de Ações?] Não --- ela não gasta
          Ação nenhuma. Mas continua limitada a uma por turno.
\end{description}

\subsection{Sobre Preparação (PP)}

\begin{description}
    \item[O Mestre pode dizer não a uma Preparação?] Só se ela sair do Escopo ou do Domínio.
          Dentro dos dois, ele não nega --- anexa uma complicação. Essa é a troca inteira.
    \item[Duas árvores de Utilidade me dão duas reservas?] Uma reserva só. Use o maior
          atributo-chave entre elas e some $+1$ por patamar de 3º ou superior em qualquer uma.
    \item[Posso preparar algo no meio de um combate?] Pode. O fato é sempre passado --- você
          está revelando, não fazendo. Mas tem que caber no que você teve tempo e motivo de
          fazer antes da cena começar.
\end{description}

\subsection{Sobre a Guilda e Combinações Novas}

\begin{description}
    \item[Meu Rank de Aventureiro sobe conforme eu gasto PA?] Não. É decisão do Mestre sobre
          feitos públicos (Cap. 5, §2) --- PA gasto não move esse marcador. O palpite que a
          ficha digital mostra é conveniência de interface, não regra.
    \item[Posso combinar magia com técnica de Corpo ou Utilidade?] Pode --- desde Rank Avançado
          nas \emph{duas} árvores envolvidas, pagando o custo de cada lado.
\end{description}

% ---------------------------------------------------------------------------
\chapter{Viagem entre Continentes}

O Mundo de Seis Faces tem seis continentes, e cruzar de um pro outro nunca é rápido nem
barato. Este apêndice dá ao Mestre uma régua rápida pra travessias longas, reaproveitando os
blocos de 1 semana do Downtime.

\begin{tabelalivro}{L{3.6cm} L{3.8cm} C{2.2cm} L{4cm}}{Rotas e Tempo de Travessia}
\cabtabela{\textbf{Rota} & \textbf{Meio} & \textbf{Tempo} & \textbf{Risco}}
Central $\leftrightarrow$ Millis & Navio de linha, porto grande & 1 bloco & Baixo --- a rota comercial mais movimentada do mundo. \\
Central $\leftrightarrow$ Begaritt & Caravana pelo deserto, ou navio pela costa & 2 blocos por terra; 1 por mar & Médio --- por terra enfrenta Clima Extremo o trajeto inteiro. \\
Millis $\leftrightarrow$ Cont. Demônio & Travessia do Estreito & 2 blocos & Alto --- águas raramente patrulhadas, sem tratado de livre passagem. \\
Central/Millis $\leftrightarrow$ Cont. Divino & Só por convite ou peregrinação registrada & 3 blocos & Baixo em trânsito, altíssimo em acesso. \\
Qualquer rota $\leftrightarrow$ Cont. Demônio por terra & Não existe. & --- & O Continente Demônio é isolado por água em todas as direções conhecidas. \\
\end{tabelalivro}

\begin{tabelalivro}{L{3.4cm} L{3.4cm} C{2.8cm} L{4cm}}{Perigo por Região}
\cabtabela{\textbf{Região} & \textbf{Teste de Clima} & \textbf{Encontro / semana} & \textbf{Nota}}
Grande Floresta (Millis) & Nenhum --- clima ameno & Alta (1d6: 1--2) & Perder-se é o perigo real, não o combate. \\
Deserto de Begaritt & CD 14 ao meio-dia & Média (1d6: 1) & Sem ração de sobra, a Sede sozinha mata uma caravana despreparada. \\
Continente Demônio & CD 8 a 18, conforme a sub-região & Alta (1d6: 1--3) & O risco não é o clima, é não ter a quem recorrer. \\
Mar aberto & Nenhum, exceto tempestade & Baixa (1d10: 1) & O maior risco é o navio, não o grupo. \\
\end{tabelalivro}

\begin{nota}{Não é sobre rolar toda semana}
    Estas tabelas existem pra resolver uma travessia em trinta segundos quando ela não é o foco
    da sessão. Se a travessia \emph{é} o foco, ignore a tabela e narre cena a cena.
\end{nota}

% ---------------------------------------------------------------------------
\chapter{Bestiário --- Criaturas por Patamar}

Em vez de um manual de monstros exaustivo, um molde por patamar calibrado com a curva que já
existe no livro.

\begin{tabelalivro}{L{3.2cm} C{1.4cm} C{1.2cm} C{2cm} C{2.4cm} C{1.8cm} C{2cm}}{}
\cabtabela{\textbf{Patamar} & \textbf{PV} & \textbf{CA} & \textbf{B. Ataque} & \textbf{Dano/turno} & \textbf{CD resist.} & \textbf{B. Resist.}}
1º --- Comum & 20 & 12 & $+3$ & \textasciitilde10 & 11 & $+2$ \\
2º --- Perigosa & 45 & 14 & $+4$ & \textasciitilde20 & 13 & $+2$ \\
3º --- Ameaça & 90 & 16 & $+6$ & \textasciitilde35 & 15 & $+3$ \\
4º --- Elite & 150 & 18 & $+8$ & \textasciitilde55 & 17 & $+4$ \\
5º --- Terror & 220 & 20 & $+10$ & \textasciitilde80 & 19 & $+5$ \\
6º --- Lenda & 320 & 22 & $+12$ & \textasciitilde120 & 21 & $+6$ \\
\end{tabelalivro}

\begin{nota}{As duas colunas que mudaram, e por quê}
    \textbf{CA} subia $+1$ por patamar, e o bônus de ataque de um personagem sobe $+1$ de Rank
    \emph{mais} o crescimento do atributo. O resultado era uma chance de acerto congelada:
    70\% com atributo 4, 90\% com atributo 8 --- a mesma coisa no 1º e no 6º patamar. Agora ela
    sobe $+2$ por patamar, e o personagem maximizado sai de 90\% pra 70\% ao longo da campanha.

    \medskip
    \textbf{Bônus de Resistência} não existia. A tabela dizia tudo de que a criatura precisava
    pra \emph{atacar}, e nada pra quando ela \emph{resiste} --- mas metade das habilidades deste
    livro pede um teste do alvo. O valor é metade do Bônus de Ataque, arredondado pra cima.
\end{nota}

\begin{aviso}{Esta tabela é a referência de calibragem do livro inteiro}
    Toda auditoria de balanceamento deste sistema mede as árvores contra estes números: dano
    por turno da criatura contra os PV do personagem, e dano por turno do personagem contra os
    PV da criatura. Uma habilidade que entregue muito acima da linha do patamar dela está fora
    da curva, mesmo que pareça divertida.
\end{aviso}
